\section{Proposed Methodology}
\label{sec:methodology}

\subsection{System Model}
We consider an edge computing network composed of $N$ heterogeneous edge servers interconnected by communication links with varying latency and bandwidth. Each server $v \in V$ is defined by a dynamic state vector $\mathbf{x}_v(t)$ that includes:
\begin{itemize}
	\item CPU utilization and available processing capacity,
	\item Memory availability and active allocation,
	\item Uplink and downlink bandwidth,
	\item Historical workload statistics,
	\item Semantic service affinity scores.
\end{itemize}

The edge fabric is modeled as a weighted graph $G = (V, E)$ where edge weights encode communication costs (latency and bandwidth). To improve scalability and reduce synchronization overhead, the network is organized into $K$ semantic-topological clusters based on geographical proximity and service similarity. Each cluster has a local orchestrator responsible for task placement before participating in periodic global policy aggregation.

\subsection{Semantic Embedding}
For each incoming task $i$, high-level descriptors such as service intent and QoS constraints are embedded into a $d$-dimensional semantic space:
\begin{equation}
	\mathbf{s}_i = f_{\text{embed}}(\text{desc}_i) \in \mathbb{R}^d,
\end{equation}
where $f_{\text{embed}}(\cdot)$ is a Sentence-BERT encoder~\cite{Reimers2019_SentenceBERT} fine-tuned on domain-specific service descriptions. Each edge node $v$ maintains a semantic profile vector $\mathbf{p}_v \in \mathbb{R}^d$ that reflects its service capabilities and historical workload characteristics.

Semantic affinity between task $i$ and node $v$ is computed using cosine similarity:
\begin{equation}
	\text{sim}(i,v) = \frac{\mathbf{s}_i \cdot \mathbf{p}_v}{\|\mathbf{s}_i\|_2 \|\mathbf{p}_v\|_2}.
\end{equation}

The semantic score is fused with numerical resource metrics to form a unified state representation. This semantic layer supports intent-aware orchestration and maintains high semantic fidelity in task assignments.

\subsection{Topology Encoding with Graph Convolutional Networks}
To capture spatial dependencies and inter-node relationships among edge servers, each cluster is represented as a graph $G_k = (V_k, E_k)$. Each node $v \in V_k$ has an initial feature vector $\mathbf{h}_v^{(0)} \in \mathbb{R}^{d_0}$ formed by concatenating semantic affinity scores, resource availability metrics, and historical load profiles.

A two-layer Graph Convolutional Network (GCN)~\cite{Kipf2017_GCN} performs message passing across the cluster topology:
\begin{equation}
	\mathbf{h}_v^{(l+1)} = \sigma\!\left(\sum_{u \in \mathcal{N}(v) \cup \{v\}} \frac{1}{\sqrt{\tilde{d}_u \tilde{d}_v}} \mathbf{W}^{(l)} \mathbf{h}_u^{(l)}\right),
\end{equation}
where:
\begin{itemize}
	\item $\mathcal{N}(v)$ is the set of neighbors of node $v$,
	\item $\tilde{d}_v = d_v + 1$ is the augmented degree accounting for self-loops,
	\item $\mathbf{W}^{(l)} \in \mathbb{R}^{d_l \times d_{l+1}}$ is a learnable weight matrix at layer $l$,
	\item $\sigma(\cdot) = \text{ReLU}(\cdot)$ is the activation function.
\end{itemize}

The final node embedding $\mathbf{h}_v^{(2)} \in \mathbb{R}^{d_h}$ (where $d_h = 64$) encodes both local resource states and multi-hop structural context from neighboring servers. These topology-aware embeddings guide orchestration decisions and reduce unnecessary service migrations by accounting for network proximity.

\subsection{Hierarchical Federated Proximal Policy Optimization}
Each cluster $k$ learns its own orchestration policy using a Proximal Policy Optimization (PPO) agent~\cite{Schulman2017_PPO} with a stochastic policy $\pi_{\theta_k}(a|s)$ that maps states (comprising GNN embeddings and resource metrics) to placement actions. The policy is optimized using the clipped surrogate objective:
\begin{equation}
	L^{\text{CLIP}}(\theta) = \mathbb{E}_t\!\left[\min\!\big(r_t(\theta)\hat{A}_t, \ \text{clip}(r_t(\theta), 1-\epsilon, 1+\epsilon)\hat{A}_t\big)\right],
\end{equation}
where $r_t(\theta) = \frac{\pi_\theta(a_t|s_t)}{\pi_{\theta_{\text{old}}}(a_t|s_t)}$ is the importance sampling ratio, $\hat{A}_t$ is the advantage estimate, and $\epsilon = 0.2$ is the clipping parameter that constrains policy updates.

The advantage function is computed using temporal difference (TD) estimation:
\begin{equation}
	\hat{A}_t = R_t + \gamma V(s_{t+1}) - V(s_t),
\end{equation}
where $V(s_t)$ is the critic's state value estimate and $\gamma = 0.99$ is the discount factor.

After every $K_1 = 10$ local training epochs, cluster-level policy parameters are synchronized within each cluster via intra-cluster federated averaging. After every $K_2 = 50$ epochs, all cluster policies are aggregated at the global coordinator:
\begin{equation}
	\theta_{\text{global}} = \frac{\sum_{k=1}^C n_k \theta_k}{\sum_{k=1}^C n_k},
\end{equation}
where $n_k$ denotes the number of training trajectories collected from cluster $k$ during the current round, and $C = 10$ is the total number of clusters. The updated global policy is then redistributed to all clusters for the next training iteration.

This hierarchical federated learning structure balances training stability with communication efficiency, enabling FedSemGNN to achieve ultra-low orchestration latency while maintaining high reward and semantic fidelity.

\begin{figure}[!t]
	\centering
	\includegraphics[width=0.95\linewidth]{Flowchart2.pdf}
	\caption{PPO-based orchestration workflow in FedSemGNN showing state encoding via GNN, action selection, reward evaluation, and policy update through hierarchical federated aggregation.}
	\label{fig:ppo_flowchart}
\end{figure}

\subsection{Reward Function}
The reward signal at timestep $t$ is designed to jointly optimize energy efficiency and response latency:
\begin{equation}
	R_t^{\text{base}} = \max\!\left(0, R_{\text{base}} - \left(P_t + w_l \cdot L_t\right)\right),
	\label{eq:base_reward}
\end{equation}
where:
\begin{itemize}
	\item $R_{\text{base}} = 10{,}000$ is the maximum achievable reward scale,
	\item $P_t$ (W) is the aggregate power consumption across all edge servers at timestep $t$,
	\item $L_t$ (ms) is the mean end-to-end service latency,
	\item $w_l = 0.01$ is the latency penalty weight that balances energy and latency objectives.
\end{itemize}

To reduce training instability caused by high-variance reward signals, we apply exponential moving average (EMA) smoothing:
\begin{equation}
	R_t = \alpha_{\text{smooth}} \cdot R_t^{\text{base}} + (1 - \alpha_{\text{smooth}}) \cdot R_{t-1},
	\label{eq:smoothed_reward}
\end{equation}
where $\alpha_{\text{smooth}} = 0.2$ controls the smoothing strength. This formulation inherently penalizes power consumption and latency while rewarding efficient orchestration decisions, eliminating the need for explicit multi-objective weighting.

Additionally, semantic fidelity is implicitly encouraged through the state representation: tasks with high semantic similarity to their assigned servers naturally achieve lower latency and better resource utilization, which translates to higher rewards. This design ensures that FedSemGNN maintains semantic alignment without requiring explicit fidelity terms in the reward function.

\subsection{Training and Synchronization Protocol}
Training proceeds in synchronous federated rounds, each consisting of $T = 500$ environment steps. Each round comprises four stages:
\begin{enumerate}
	\item \textbf{Trajectory Collection}: Edge nodes execute their current local policies and collect state-action-reward-next-state tuples into a prioritized experience replay buffer with capacity 20,000.
	\item \textbf{Local PPO Optimization}: Each cluster orchestrator performs $E = 4$ epochs of PPO updates using mini-batches of size $B = 64$ sampled from the replay buffer, optimizing the clipped surrogate objective with learning rate $\alpha = 0.001$.
	\item \textbf{Hierarchical Policy Aggregation}: 
	\begin{itemize}
		\item Every $K_1 = 10$ epochs: Intra-cluster aggregation synchronizes policies within each cluster.
		\item Every $K_2 = 50$ epochs: Inter-cluster aggregation computes the global policy $\theta_{\text{global}}$ via weighted averaging across all $C = 10$ clusters.
	\end{itemize}
	\item \textbf{Global Distribution}: The updated global policy is broadcast to all cluster orchestrators for the next training round.
\end{enumerate}

Gradient clipping with maximum norm 0.5 and differential privacy noise ($\sigma_{\text{DP}} = 1.0$) are applied during aggregation to ensure training stability and privacy preservation. This protocol enables FedSemGNN to maintain ultra-low orchestration latency (\textbf{0.36 ms per step})—over 315$\times$ faster than FlatFedPPO (\textbf{114.25 ms}) and 222$\times$ faster than HierFedPPO (\textbf{80.71 ms}).

\subsection{Implementation Details}
FedSemGNN is implemented using PyTorch 2.0 and PyTorch Geometric 2.3 for GCN modules. The PPO agent employs a two-layer fully connected actor-critic architecture with $h_{\text{actor}} = h_{\text{critic}} = 16$ hidden units per layer, ReLU activations, and dropout rate $p_{\text{drop}} = 0.5$ for regularization. The GCN encoder uses two graph convolutional layers ($L = 2$) with $d_h = 64$ dimensional hidden embeddings and symmetric normalization $\tilde{D}^{-\frac{1}{2}} \tilde{A} \tilde{D}^{-\frac{1}{2}}$.

Semantic embeddings are generated using the pre-trained \texttt{all-MiniLM-L6-v2} Sentence-BERT model~\cite{Reimers2019_SentenceBERT}, which produces 384-dimensional vectors that are then projected to the optimized $d_s = 128$ dimensional semantic space for computational efficiency. Online semantic learning with elastic weight consolidation ($\lambda_{\text{ewc}} = 0.4$) enables real-time adaptation to evolving service requirements at learning rate $\alpha_{\text{online}} = 0.0001$.

All experiments are conducted using EdgeSimPy~\cite{edgesimpy}, a realistic edge computing simulator with detailed network topology modeling and physical-layer constraints. The base simulation environment comprises 6 heterogeneous edge servers with varying computational capacities (4-16 CPU cores, 8-32 GB RAM), 500 mobile users with dynamic mobility patterns following random waypoint models, and heterogeneous service requests representing diverse QoS requirements (latency-sensitive, compute-intensive, bandwidth-heavy). The network topology follows a mesh connectivity pattern generated using NetworkX.

The simulation executes for 500 timesteps with tick duration $\Delta t = 1$ second. Each configuration is evaluated across 5 independent trials with different random seeds, and all reported results represent mean values $\pm$ standard deviation to ensure statistical significance.

Table~\ref{tab:hyperparameters} summarizes all hyperparameters and configuration settings used in FedSemGNN.

% [Hyperparameters table included here as provided]

\subsection{Scalability Evaluation Methodology}
To evaluate scalability at extreme scales (1,000, 5,000, and 10,000 nodes), we employ a mathematical extrapolation methodology validated in prior distributed systems research~\cite{Bonawitz2019TowardsFL, Li2020_FederatedOptimization}. We fit power-law and logarithmic scaling curves to the base simulation results for key performance metrics (latency, power consumption, communication overhead, reward), then extrapolate these trends to larger node counts with confidence interval estimation.

This approach enables rigorous scalability analysis without the prohibitive computational cost of simulating massive topologies directly (which would require weeks of computation time), while maintaining statistical rigor through validated scaling models from the distributed systems literature. The extrapolation assumes:
\begin{itemize}
	\item Communication overhead scales logarithmically with hierarchical clustering,
	\item Power consumption scales linearly with active nodes,
	\item Orchestration latency follows sub-linear growth due to parallelization,
	\item Reward convergence maintains similar trends across scales.
\end{itemize}

For large-scale configurations, hyperparameters are adjusted to maintain stability: learning rate is reduced to $\alpha = 10^{-4}$, batch size decreased to $B = 32$, clipping parameter tightened to $\epsilon = 0.15$, and cluster sizes adapted dynamically between 20-200 nodes (see Table~\ref{tab:hyperparameters} notes).

Table~\ref{tab:performance_comparison} summarizes the comparative performance of FedSemGNN against state-of-the-art baseline algorithms across all evaluation metrics.
